\clearpage
\VerdierPhysicalPageAnchor{0162}%
\begingroup\small\setlength{\parskip}{0pt}\setlength{\emergencystretch}{1.5em}
\VerdierSemanticAnchor{ch4}
\section*{4. 谱对象。}

我们重新采用 [1] 第 XV 章 \S 7 对谱序列的论述，并将其改写以适用于
三角范畴。

\VerdierSemanticAnchor{ch4-4-1}
\subsection*{4.1. 谱对象的定义。}

\VerdierSemanticAnchor{ch4-4-1-1}
\medskip\noindent\textbf{4.1.1.}\quad
设 $\mathcal J$ 为一个范畴，并且对每个 $n\in\mathbb N$，以
$\mathcal{F}\ell_n(\mathcal J)$ 表示由 $n$ 个 $\mathcal J$ 中可复合态射
组成的列的范畴（第 I 章，3.3.6）。我们将使用三个记为
$s_i\colon\mathcal{F}\ell_2(\mathcal J)\to\mathcal{F}\ell_1(\mathcal J)$，
$1\leq i\leq3$ 的函子：
\VerdierSemanticAnchor{ch4-4-1-1-1}
\[
\left\{
\begin{aligned}
&\bigl(\mathord\bullet\xrightarrow f\mathord\bullet\xrightarrow g\mathord\bullet\bigr)
   \in\operatorname{Ob}\mathcal{F}\ell_2(\mathcal J),\\
&s_1\bigl(\mathord\bullet\xrightarrow f\mathord\bullet\xrightarrow g\mathord\bullet\bigr)
   =\bigl(\mathord\bullet\xrightarrow f\mathord\bullet\bigr),\\
&s_2\bigl(\mathord\bullet\xrightarrow f\mathord\bullet\xrightarrow g\mathord\bullet\bigr)
   =\bigl(\mathord\bullet\xrightarrow{gf}\mathord\bullet\bigr),\\
&s_3\bigl(\mathord\bullet\xrightarrow f\mathord\bullet\xrightarrow g\mathord\bullet\bigr)
   =\bigl(\mathord\bullet\xrightarrow g\mathord\bullet\bigr).
\end{aligned}
\right.
\tag{4.1.1.1}
\]
以及两个函子态射：
\VerdierSemanticAnchor{ch4-4-1-1-2}
\[
\left\{
\begin{aligned}
&u_1\colon s_1\longrightarrow s_2,
&&u_1(f,g)=
\begin{tikzcd}[ampersand replacement=\&,column sep=2.8em,row sep=1.9em,baseline=(current bounding box.center)]
{}\arrow[r,"f"]\arrow[d,"\operatorname{id}"']\&{}\arrow[d,"g"]\\
{}\arrow[r,"gf"']\&{}
\end{tikzcd},\\[2pt]
&u_2\colon s_2\longrightarrow s_3,
&&u_2(f,g)=
\begin{tikzcd}[ampersand replacement=\&,column sep=2.8em,row sep=1.9em,baseline=(current bounding box.center)]
{}\arrow[r,"gf"]\arrow[d,"f"']\&{}\arrow[d,"\operatorname{id}"]\\
{}\arrow[r,"g"']\&{}
\end{tikzcd}.
\end{aligned}
\right.
\tag{4.1.1.2}
\]

\VerdierSemanticAnchor{ch4-4-1-2}
\medskip\noindent\textbf{4.1.2.}\quad
设 $\mathcal D$ 为三角范畴，$\mathcal J$ 为一个范畴。
一个\emph{$\mathcal J$ 型且取值于 $\mathcal D$ 的谱对象}，是由以下两个
函子组成的对象：

a) 一个函子
$X\colon\mathcal{F}\ell_1(\mathcal J)\longrightarrow\mathcal D$：
\[
\bigl(\mathord\bullet\xrightarrow f\mathord\bullet\bigr)\longmapsto X_f.
\]
\endgroup
